% ---------------------------------------------------------------------------
% Shared preamble for "Memory in Large Language Models" (pdflatex).
% Chapters must use only the packages and macros defined here.
% ---------------------------------------------------------------------------
\usepackage[T1]{fontenc}
\usepackage[utf8]{inputenc}
\usepackage{lmodern}
\usepackage[a4paper,inner=1.1in,outer=0.9in,top=1in,bottom=1.1in,headheight=15pt]{geometry}
\usepackage{microtype}
\usepackage{textcomp}

% Mathematics
\usepackage{amsmath,amssymb,amsthm,mathtools}
\usepackage{bm}

% Tables, lists, figures
\usepackage{graphicx}
\usepackage{booktabs,tabularx,longtable,array,multirow}
\usepackage{enumitem}
\usepackage{caption}
\usepackage{xcolor}
\usepackage{tikz}
\usetikzlibrary{arrows.meta,positioning,calc,fit,shapes.geometric,decorations.pathreplacing,backgrounds}

% Code and algorithms
\usepackage{listings}
\usepackage{algorithm}
\usepackage[noend]{algpseudocode}

% Boxes
\usepackage[most]{tcolorbox}

% Headers
\usepackage{fancyhdr}

% References
\usepackage[numbers,sort&compress]{natbib}
\usepackage{url}
\usepackage[colorlinks=true,linkcolor=blue!50!black,citecolor=green!40!black,urlcolor=magenta!60!black,bookmarksnumbered=true]{hyperref}
\usepackage[capitalise,noabbrev]{cleveref}

% ---------------------------------------------------------------------------
% Page style
% ---------------------------------------------------------------------------
\pagestyle{fancy}
\fancyhf{}
\fancyhead[LE,RO]{\thepage}
\fancyhead[RE]{\small\nouppercase{\leftmark}}
\fancyhead[LO]{\small\nouppercase{\rightmark}}
\renewcommand{\headrulewidth}{0.4pt}
\fancypagestyle{plain}{\fancyhf{}\fancyfoot[C]{\thepage}\renewcommand{\headrulewidth}{0pt}}

\setlength{\parskip}{0.3em}
\setlist{itemsep=0.15em,topsep=0.3em}
\captionsetup{font=small,labelfont=bf}
\numberwithin{equation}{chapter}
\renewcommand{\arraystretch}{1.15}

% ---------------------------------------------------------------------------
% Theorem-like environments (numbered per chapter)
% ---------------------------------------------------------------------------
\theoremstyle{definition}
\newtheorem{definition}{Definition}[chapter]
\newtheorem{example}[definition]{Example}
\newtheorem{exercise}{Exercise}[chapter]
\theoremstyle{plain}
\newtheorem{theorem}[definition]{Theorem}
\newtheorem{proposition}[definition]{Proposition}
\newtheorem{lemma}[definition]{Lemma}
\newtheorem{corollary}[definition]{Corollary}
\theoremstyle{remark}
\newtheorem{remark}[definition]{Remark}
\crefname{definition}{Definition}{Definitions}
\crefname{example}{Example}{Examples}
\crefname{exercise}{Exercise}{Exercises}
\crefname{proposition}{Proposition}{Propositions}
\crefname{lemma}{Lemma}{Lemmas}
\crefname{corollary}{Corollary}{Corollaries}
\crefname{remark}{Remark}{Remarks}
\crefname{lstlisting}{Listing}{Listings}
\crefname{algorithm}{Algorithm}{Algorithms}

% ---------------------------------------------------------------------------
% Callout boxes. Optional argument passes tcolorbox keys, e.g.
%   \begin{implnote}[title={In the code: fla}] ... \end{implnote}
% ---------------------------------------------------------------------------
\tcbset{membox/.style={enhanced,breakable,boxrule=0.6pt,arc=2pt,left=6pt,right=6pt,top=4pt,bottom=4pt,fonttitle=\bfseries\small,before skip=8pt,after skip=8pt}}
\newtcolorbox{keyidea}[1][]{membox,colback=blue!3,colframe=blue!45!black,title={Key idea},#1}
\newtcolorbox{taxonomy}[1][]{membox,colback=violet!3,colframe=violet!50!black,title={Where it sits in the taxonomy},#1}
\newtcolorbox{implnote}[1][]{membox,colback=teal!3,colframe=teal!50!black,title={In the code},#1}
\newtcolorbox{results}[1][]{membox,colback=orange!4,colframe=orange!60!black,title={Reported results},#1}
\newtcolorbox{pitfall}[1][]{membox,colback=red!3,colframe=red!50!black,title={Pitfall},#1}
\newtcolorbox{summarybox}[1][]{membox,colback=gray!5,colframe=gray!50!black,title={Chapter summary},#1}

% ---------------------------------------------------------------------------
% Code listings (ASCII only inside listings)
% ---------------------------------------------------------------------------
\definecolor{codebg}{gray}{0.97}
\definecolor{codekw}{RGB}{0,70,150}
\definecolor{codecomment}{RGB}{90,110,90}
\definecolor{codestr}{RGB}{160,40,40}
\lstdefinestyle{book}{
  language=Python,
  basicstyle=\ttfamily\footnotesize,
  keywordstyle=\color{codekw}\bfseries,
  commentstyle=\color{codecomment}\itshape,
  stringstyle=\color{codestr},
  backgroundcolor=\color{codebg},
  frame=single, framerule=0.3pt, rulecolor=\color{gray!50},
  numbers=left, numberstyle=\tiny\color{gray}, numbersep=6pt,
  breaklines=true, breakatwhitespace=false, postbreak=\mbox{\textcolor{gray}{$\hookrightarrow$}\space},
  showstringspaces=false, columns=fullflexible, keepspaces=true, upquote=true,
  tabsize=4, captionpos=t, xleftmargin=14pt, xrightmargin=2pt,
  morekeywords={self,None,True,False,torch,nn,einsum}
}
\lstset{style=book}

% Pseudocode tweaks
\algrenewcommand\algorithmicrequire{\textbf{Input:}}
\algrenewcommand\algorithmicensure{\textbf{Output:}}
\algnewcommand{\LineComment}[1]{\State \(\triangleright\) \textit{#1}}

% ---------------------------------------------------------------------------
% Notation
% ---------------------------------------------------------------------------
\newcommand{\R}{\mathbb{R}}
\newcommand{\C}{\mathbb{C}}
\newcommand{\N}{\mathbb{N}}
\newcommand{\E}{\mathbb{E}}
\renewcommand{\vec}[1]{\bm{#1}}
% vectors (bold lowercase)
\newcommand{\va}{\bm{a}}
\newcommand{\vb}{\bm{b}}
\newcommand{\ve}{\bm{e}}
\newcommand{\vg}{\bm{g}}
\newcommand{\vh}{\bm{h}}
\newcommand{\vk}{\bm{k}}
\newcommand{\vm}{\bm{m}}
\newcommand{\vo}{\bm{o}}
\newcommand{\vq}{\bm{q}}
\newcommand{\vr}{\bm{r}}
\newcommand{\vs}{\bm{s}}
\newcommand{\vu}{\bm{u}}
\newcommand{\vv}{\bm{v}}
\newcommand{\vw}{\bm{w}}
\newcommand{\vx}{\bm{x}}
\newcommand{\vy}{\bm{y}}
\newcommand{\vz}{\bm{z}}
\newcommand{\valpha}{\bm{\alpha}}
\newcommand{\vbeta}{\bm{\beta}}
\newcommand{\vtheta}{\bm{\theta}}
% matrices (bold uppercase)
\newcommand{\mA}{\bm{A}}
\newcommand{\mB}{\bm{B}}
\newcommand{\mC}{\bm{C}}
\newcommand{\mD}{\bm{D}}
\newcommand{\mG}{\bm{G}}
\newcommand{\mH}{\bm{H}}
\newcommand{\mI}{\bm{I}}
\newcommand{\mK}{\bm{K}}
\newcommand{\mM}{\bm{M}}
\newcommand{\mO}{\bm{O}}
\newcommand{\mP}{\bm{P}}
\newcommand{\mQ}{\bm{Q}}
\newcommand{\mS}{\bm{S}}
\newcommand{\mT}{\bm{T}}
\newcommand{\mU}{\bm{U}}
\newcommand{\mV}{\bm{V}}
\newcommand{\mW}{\bm{W}}
\newcommand{\mX}{\bm{X}}
\newcommand{\mY}{\bm{Y}}
\newcommand{\mZ}{\bm{Z}}
\newcommand{\mLambda}{\bm{\Lambda}}
\newcommand{\mSigma}{\bm{\Sigma}}
% operators
\DeclareMathOperator{\softmax}{softmax}
\DeclareMathOperator{\diag}{diag}
\DeclareMathOperator{\tr}{tr}
\DeclareMathOperator{\RMSNorm}{RMSNorm}
\DeclareMathOperator{\LayerNorm}{LayerNorm}
\DeclareMathOperator{\SiLU}{SiLU}
\DeclareMathOperator{\ReLU}{ReLU}
\DeclareMathOperator{\TopK}{TopK}
\DeclareMathOperator{\rank}{rank}
\DeclareMathOperator*{\argmin}{arg\,min}
\DeclareMathOperator*{\argmax}{arg\,max}
\newcommand{\T}{^{\top}}
\newcommand{\norm}[1]{\left\lVert #1\right\rVert}
\newcommand{\abs}[1]{\left\lvert #1\right\rvert}
\newcommand{\inner}[2]{\left\langle #1,\, #2\right\rangle}
\newcommand{\bigO}{\mathcal{O}}
\newcommand{\loss}{\mathcal{L}}
\newcommand{\sigmoid}{\sigma}
\newcommand{\had}{\odot}
\newcommand{\defeq}{\coloneqq}

% ---------------------------------------------------------------------------
% Code and repository references (do not use inside section titles or captions;
% there, write \texttt{...} with \_ escapes instead)
% ---------------------------------------------------------------------------
\newcommand{\code}[1]{\nolinkurl{#1}}        % inline identifier, e.g. \code{chunk_gated_delta_rule}
\newcommand{\repofile}[1]{\nolinkurl{#1}}    % path relative to repo root, e.g. \repofile{repos/NVlabs__GatedDeltaNet/README.md}
\newcommand{\repo}[1]{\nolinkurl{repos/#1}}  % mirror directory, e.g. \repo{NVlabs__GatedDeltaNet}
\newcommand{\gh}[1]{\href{https://github.com/#1}{\nolinkurl{#1}}} % upstream GitHub repo, e.g. \gh{NVlabs/GatedDeltaNet}
\newcommand{\fla}{\textsf{fla}}
\newcommand{\arxiv}[1]{\href{https://arxiv.org/abs/#1}{arXiv:#1}}
